\section{14.2 word2vec: 단어를 벡터로}
\begin{frame}{임베딩: PCA와는 다른 방향의 압축}
  \begin{columns}[T]
    \begin{column}{0.48\textwidth}
      \kb{PCA}
      \begin{itemize}
        \item ``이미 \textbf{벡터인} 데이터''를 회전·압축.
        \item 기준: \textbf{분산 보존} (데이터 내적).
      \end{itemize}
      \vskip0.6em
      \kb{임베딩}(embedding)
      \begin{itemize}
        \item 원래 벡터가 \textbf{아니었던} 것(단어, 그래프 노드)을
          \textbf{학습}으로 저차원 벡터로.
        \item 기준: \textbf{예측 성능} (문맥 예측).
      \end{itemize}
    \end{column}
    \begin{column}{0.5\textwidth}
      \begin{block}{같은 목표}
        고차원(또는 벡터가 아닌) 원본을, \textbf{유용한 구조를 보존한
        채} 저차원 벡터로 압축.
      \end{block}
      \vskip0.4em
      \kq{함께 자주 나타나는 것들은 가까운 벡터로 모인다.}
      \; --- 이 장 전체의 한 줄 원칙.
    \end{column}
  \end{columns}
  \vskip0.6em
  {\footnotesize 2013년 미콜로프(Tomas Mikolov) 등이 제안한
  \kb{word2vec} (Mikolov et al., 2013)은 ``단어의 의미는 주변 단어로
  결정된다''는 \kb{분포 가설}(distributional hypothesis)을 학습으로
  구현한 것.}
\end{frame}
